\section{Experimental Setup}
\label{sec:experimental-setup}
% TODO: Environments (Passive T-Maze, Noisy T-Maze, Associative
% Corridor; Car Flag and Memory Cards as no-loss checks), baselines
% (DTQN, DRQN), hyperparameters, seeds, metrics (success rate,
% 95% CI, wall clock, peak GPU memory).
% refs: \cite{esslinger2022dtqn,hausknecht2015drqn,mnih2015dqn}

\paragraph{Evaluation protocol and criteria.}
Each configuration is trained with five seeds and evaluated greedily on 100
episodes; we report the interquartile mean (IQM), mean, median and optimality gap
of the final success rate with 95\% stratified bootstrap confidence intervals and
the standard deviation over runs, the final return min-max normalized per task,
and the probability of improvement $P(X > Y)$ \cite{agarwal2021rliable}.
A configuration is \emph{beyond the window} when the cue is older than the
context, i.e.\ when the context holds fewer than $L + 1$ observations for
corridor length $L$, and \emph{within the window} otherwise. Beyond the window,
the memory succeeds when DTQN-kNN reaches an IQM success rate of at least $0.9$
and the 95\% confidence interval of $P(\text{DTQN-kNN} > \text{DTQN})$ lies
above $0.5$. Within the window, DTQN-kNN ties with DTQN (non-inferiority) when
that interval is not entirely below $0.5$ and its IQM is at most $0.05$ below
that of DTQN. A criterion that is not met is reported as such.
